\chapter{Cell Attachments}
\label{ch:viii05-cell-attachments}

Cell attachment turns local disks into global spaces.  Starting from \(X\), an
\(n\)-disk is glued along its boundary by a map
\[
\varphi:S^{n-1}\to X.
\]
The quotient remembers the old space, the new cell dimension, and the attaching
map.  This chapter develops the full adjunction-space language and the major
problem families mapped to cell attachments, while reserving the theorem about
homotopic attaching maps for VIII/08.

\section{Adjunction spaces}
\begin{definition}[Attaching an \(n\)-cell]\label{def:viii05-attach}
\[
X\cup_\varphi D^n=(X\sqcup D^n)/\!\sim,
\]
where \(u\sim\varphi(u)\) for \(u\in S^{n-1}=\partial D^n\).
\end{definition}

\section{Characteristic map and open cell}
Let \(q:X\sqcup D^n\to X\cup_\varphi D^n\) be the quotient map.

\begin{definition}[Characteristic map]\label{def:viii05-characteristic}
The characteristic map is
\[
\chi=q|_{D^n}:D^n\to X\cup_\varphi D^n.
\]
Its restriction to \(\mathring D^n\) is a homeomorphism onto the open
\(n\)-cell \(e^n\).
\end{definition}

\begin{warning}[Boundary identifications]\label{warn:viii05-boundary}
The characteristic map need not embed the closed disk.  Distinct boundary
points can be identified by the attaching map.
\end{warning}

\section{Universal property}
\begin{proposition}[Gluing criterion]\label{prop:viii05-universal}
To define \(F:X\cup_\varphi D^n\to Y\) is equivalent to giving continuous maps
\(f:X\to Y\) and \(g:D^n\to Y\) satisfying
\[
g|_{S^{n-1}}=f\circ\varphi.
\]
\end{proposition}
This is the quotient, equivalently pushout, universal property.

\section{Zero- and one-cells}
A \(0\)-cell is a point and has empty boundary, so attaching it adds a new
disjoint point.  A \(1\)-cell has boundary
\[
S^0=\{-1,+1\},
\]
so its attaching map chooses two points of the old space, possibly equal.

\begin{example}[Loop]\label{ex:viii05-loop}
Attaching both endpoints of \(D^1\) to one vertex produces \(S^1\).
\end{example}

\begin{example}[Graphs]\label{ex:viii05-graphs}
Finite graphs arise from finitely many zero-cells followed by finitely many
one-cell attachments.
\end{example}

\section{Constant attachments}
\begin{proposition}[Constant attachment]\label{prop:viii05-constant}
If \(S^{n-1}\to X\) is constant at a basepoint \(x_0\), the whole boundary of
the new disk collapses at \(x_0\).  In the standard pointed setting the result
is
\[
X\vee S^n.
\]
In particular,
\[
D^n/S^{n-1}\cong S^n.
\]
\end{proposition}

\section{Degree-\(m\) two-cell attachments}
For
\[
\varphi_m:S^1\to S^1,\qquad \varphi_m(z)=z^m,
\]
define
\[
X_m=S^1\cup_{\varphi_m}D^2.
\]

\begin{example}[Degree one]\label{ex:viii05-degree-one}
\(X_1\cong D^2\).
\end{example}

\begin{example}[Degree zero]\label{ex:viii05-degree-zero}
For a constant attaching map,
\[
X_0\cong S^1\vee S^2.
\]
\end{example}

\begin{example}[Projective plane]\label{ex:viii05-rp2}
\[
\mathbb RP^2\cong S^1\cup_{\varphi_2}D^2.
\]
\end{example}

The identical cell counts can therefore carry very different topology because
the attaching maps differ.

\section{Sphere with two cells}
\begin{proposition}[Minimal sphere structure]\label{prop:viii05-sphere}
For \(n\ge1\), \(S^n\) is obtained from one zero-cell by attaching one
\(n\)-cell with constant attaching map.
\end{proposition}

\section{Torus polygon}
The standard torus has one \(0\)-cell, two \(1\)-cells \(a,b\), and one
\(2\)-cell.  Its one-skeleton is
\[
S^1\vee S^1,
\]
and the two-cell attaches by the commutator word
\[
aba^{-1}b^{-1}.
\]

\section{Closed orientable surfaces}
A genus-\(g\) orientable surface has one \(0\)-cell, \(2g\) one-cells, and one
two-cell, with attaching word
\[
[a_1,b_1]\cdots[a_g,b_g].
\]

\section{Projective filtration}
\begin{proposition}[Projective cells]\label{prop:viii05-projective}
\(\mathbb RP^n\) is obtained from \(\mathbb RP^{n-1}\) by attaching one
\(n\)-cell.  Hence it has one cell in every dimension \(0,\dots,n\).
\end{proposition}

\section{Attaching several cells}
For a family \(\varphi_\alpha:S^{n-1}\to X\), attach all cells simultaneously:
\[
\left(X\sqcup\coprod_\alpha D^n_\alpha\right)/\!\sim,
\qquad
u\sim\varphi_\alpha(u).
\]
Distinct disk interiors remain disjoint.

\section{Successive attachment}
Dimension-by-dimension construction yields
\[
X^0\subset X^1\subset X^2\subset\cdots,
\]
where \(X^n\) is obtained from \(X^{n-1}\) by attaching \(n\)-cells.  VIII/06
formalizes these \(X^n\) as skeleta.

\section{Attaching versus characteristic maps}
The attaching map
\[
\varphi:S^{n-1}\to X
\]
describes how the boundary meets the old space.  The characteristic map
\[
\chi:D^n\to X\cup_\varphi D^n
\]
parametrizes the closure of the new cell.

\section{Pushout formulation}
Cell attachment is the pushout
\[
\begin{array}{ccc}
S^{n-1}&\xrightarrow{\varphi}&X\\
\big\downarrow&&\big\downarrow\\
D^n&\longrightarrow&X\cup_\varphi D^n.
\end{array}
\]

\section{Dependence on the attaching map}
\begin{warning}[The map matters]\label{warn:viii05-map-matters}
The number and dimensions of cells do not generally determine the resulting
space.  The attaching maps are essential data.
\end{warning}

Homotopic attaching maps often lead to homotopy-equivalent adjunction spaces,
but the rigorous theorem and its homotopy-extension argument belong to VIII/08.
Blocks explicitly mapped there are therefore not absorbed into this chapter.

\section{Mapping-cone preview}
For a map \(f:A\to X\), the mapping cone is
\[
C_f=X\cup_f CA.
\]
Since \(CS^{n-1}\cong D^n\), ordinary cell attachment is the sphere-domain
special case.  VIII/07 develops mapping cones in full.

\section{Why attachments compute topology}
Cellular homology later converts \(n\)-cell attaching maps into boundary
homomorphisms.  The degree-\(m\) family already predicts the integer \(m\) that
will appear in a cellular boundary map.

\begin{warning}[Use the quotient topology]\label{warn:viii05-quotient}
The construction is not merely a set-theoretic union.  Continuity is controlled
by the quotient topology induced by the boundary identifications.
\end{warning}


\section{Exercises}

\begin{exercise}\label{exr:viii05-01}
Write the quotient defining \(X\cup_\varphi D^n\).
\end{exercise}
\begin{hint}
Begin with the disjoint union.

% VIII-PEDAGOGY-HINT-v1.1 exr:viii05-01 append
Identify only boundary points with their prescribed images; explain which equivalence classes contain points of the disk interior.
\end{hint}
\begin{solution}
\((X\sqcup D^n)/\!\sim\), with \(u\sim\varphi(u)\) on \(S^{n-1}\).
\end{solution}

\begin{exercise}\label{exr:viii05-02}
What is the domain of an \(n\)-cell attaching map?
\end{exercise}
\begin{hint}
Use the disk boundary.

% VIII-PEDAGOGY-HINT-v1.1 exr:viii05-02 append
For \(n>0\), compare the dimensions of the disk, its interior, and its boundary; treat the zero-cell boundary as empty.
\end{hint}
\begin{solution}
\(S^{n-1}\).
\end{solution}

\begin{exercise}\label{exr:viii05-03}
Define the characteristic map.
\end{exercise}
\begin{hint}
Restrict the quotient map.

% VIII-PEDAGOGY-HINT-v1.1 exr:viii05-03 append
Start with the quotient map on \(X\sqcup D^n\), then restrict its domain to the new disk, including its boundary.
\end{hint}
\begin{solution}
\(\chi=q|_{D^n}\).
\end{solution}

\begin{exercise}\label{exr:viii05-04}
Where is \(\chi\) a homeomorphism?
\end{exercise}
\begin{hint}
Use the disk interior.

% VIII-PEDAGOGY-HINT-v1.1 exr:viii05-04 append
Check that no interior point is identified with another point; do not extend this injectivity claim to the disk boundary.
\end{hint}
\begin{solution}
On \(\mathring D^n\), onto the open cell.
\end{solution}

\begin{exercise}\label{exr:viii05-05}
State the compatibility for maps \(f:X\to Y\), \(g:D^n\to Y\).
\end{exercise}
\begin{hint}
Compare on the boundary.

% VIII-PEDAGOGY-HINT-v1.1 exr:viii05-05 append
A point \(u\) on the boundary represents both a disk point and \(\varphi(u)\) in the old space; their proposed images must agree.
\end{hint}
\begin{solution}
\(g|_{S^{n-1}}=f\circ\varphi\).
\end{solution}

\begin{exercise}\label{exr:viii05-06}
What does a zero-cell attachment do?
\end{exercise}
\begin{hint}
Its boundary is empty.

% VIII-PEDAGOGY-HINT-v1.1 exr:viii05-06 append
Since there are no boundary points to identify, use the disjoint-union topology to decide how the new point meets \(X\).
\end{hint}
\begin{solution}
It adds a new disjoint point.
\end{solution}

\begin{exercise}\label{exr:viii05-07}
What data specifies a one-cell attachment?
\end{exercise}
\begin{hint}
\(S^0\) has two points.

% VIII-PEDAGOGY-HINT-v1.1 exr:viii05-07 append
The two endpoint images may coincide or differ; compare the loop-forming and edge-joining possibilities.
\end{hint}
\begin{solution}
The images of the two endpoints.
\end{solution}

\begin{exercise}\label{exr:viii05-08}
Build \(S^1\) from one zero-cell and one one-cell.
\end{exercise}
\begin{hint}
Attach both endpoints to one vertex.

% VIII-PEDAGOGY-HINT-v1.1 exr:viii05-08 append
Parametrize the interval by \(t\mapsto e^{2\pi it}\), and check precisely which endpoint identification makes this descend.
\end{hint}
\begin{solution}
The quotient \(D^1/\partial D^1\) is \(S^1\).
\end{solution}

\begin{exercise}\label{exr:viii05-09}
How are finite graphs built cellularly?
\end{exercise}
\begin{hint}
Use dimensions 0 and 1.

% VIII-PEDAGOGY-HINT-v1.1 exr:viii05-09 append
Assign each edge an interval and record its two endpoint vertices; loops and multiple edges are allowed in this CW description.
\end{hint}
\begin{solution}
Attach finitely many one-cells to finitely many vertices.
\end{solution}

\begin{exercise}\label{exr:viii05-10}
What is \(D^n/S^{n-1}\)?
\end{exercise}
\begin{hint}
Collapse the boundary.

% VIII-PEDAGOGY-HINT-v1.1 exr:viii05-10 append
Use a disk as a compactification of its interior after collapsing the boundary; check that the collapsed class is the single point at infinity.
\end{hint}
\begin{solution}
\(S^n\).
\end{solution}

\begin{exercise}\label{exr:viii05-11}
What does a constant \(n\)-attachment give in a pointed space?
\end{exercise}
\begin{hint}
The new sphere meets at the basepoint.

% VIII-PEDAGOGY-HINT-v1.1 exr:viii05-11 append
The entire boundary is identified to the chosen basepoint, while the disk interior stays separate from the old space.
\end{hint}
\begin{solution}
\(X\vee S^n\).
\end{solution}

\begin{exercise}\label{exr:viii05-12}
Define the degree-\(m\) attaching map on \(S^1\).
\end{exercise}
\begin{hint}
Use a power map.

% VIII-PEDAGOGY-HINT-v1.1 exr:viii05-12 append
Write \(z=e^{i\theta}\) and let the output angle advance by \(m\theta\); negative integers must reverse winding direction.
\end{hint}
\begin{solution}
\(\varphi_m(z)=z^m\).
\end{solution}

\begin{exercise}\label{exr:viii05-13}
Identify \(S^1\cup_{\varphi_1}D^2\).
\end{exercise}
\begin{hint}
Boundary goes once around.

% VIII-PEDAGOGY-HINT-v1.1 exr:viii05-13 append
The old circle is identified with the new disk's actual boundary, not collapsed to a point as in the sphere construction.
\end{hint}
\begin{solution}
\(D^2\).
\end{solution}

\begin{exercise}\label{exr:viii05-14}
Identify a constant two-cell attachment to \(S^1\).
\end{exercise}
\begin{hint}
Use a wedge.

% VIII-PEDAGOGY-HINT-v1.1 exr:viii05-14 append
Keep the old circle and replace the disk with its boundary quotient; identify the single point where these two pieces meet.
\end{hint}
\begin{solution}
\(S^1\vee S^2\).
\end{solution}

\begin{exercise}\label{exr:viii05-15}
Give the cell attachment for \(\mathbb RP^2\).
\end{exercise}
\begin{hint}
Use degree two.

% VIII-PEDAGOGY-HINT-v1.1 exr:viii05-15 append
In the disk model of \(\mathbb RP^2\), opposite boundary points represent the same line; track the boundary map after identifying \(\mathbb RP^1\) with a circle.
\end{hint}
\begin{solution}
\(S^1\cup_{\varphi_2}D^2\).
\end{solution}

\begin{exercise}\label{exr:viii05-16}
Give minimal cell counts for \(S^n\).
\end{exercise}
\begin{hint}
One old point, one new cell.

% VIII-PEDAGOGY-HINT-v1.1 exr:viii05-16 append
For \(n>0\), collapse the boundary of one \(n\)-disk to the old vertex; for \(n=0\), regard \(S^0\) as two vertices.
\end{hint}
\begin{solution}
One \(0\)-cell and one \(n\)-cell.
\end{solution}

\begin{exercise}\label{exr:viii05-17}
Give the torus cell counts.
\end{exercise}
\begin{hint}
Use the square model.

% VIII-PEDAGOGY-HINT-v1.1 exr:viii05-17 append
Track the equivalence classes of the square's four vertices and its pairs of opposite edges before counting the open interior.
\end{hint}
\begin{solution}
One \(0\)-cell, two \(1\)-cells, one \(2\)-cell.
\end{solution}

\begin{exercise}\label{exr:viii05-18}
What word attaches the torus two-cell?
\end{exercise}
\begin{hint}
Use the commutator.

% VIII-PEDAGOGY-HINT-v1.1 exr:viii05-18 append
Choose orientations on the two identified edges and read the oriented square boundary once, using inverses when traversal opposes the chosen edge direction.
\end{hint}
\begin{solution}
\(aba^{-1}b^{-1}\).
\end{solution}

\begin{exercise}\label{exr:viii05-19}
Give genus-\(g\) surface cell counts.
\end{exercise}
\begin{hint}
Generalize the torus.

% VIII-PEDAGOGY-HINT-v1.1 exr:viii05-19 append
Use a \(4g\)-gon with paired sides; count the resulting vertex class, distinct edge pairs, and open polygon interior.
\end{hint}
\begin{solution}
One \(0\)-cell, \(2g\) one-cells, one two-cell.
\end{solution}

\begin{exercise}\label{exr:viii05-20}
Give its attaching word.
\end{exercise}
\begin{hint}
Multiply commutators.

% VIII-PEDAGOGY-HINT-v1.1 exr:viii05-20 append
Group the oriented boundary traversal into one four-letter block for each handle, retaining the order of the blocks.
\end{hint}
\begin{solution}
\([a_1,b_1]\cdots[a_g,b_g]\).
\end{solution}

\begin{exercise}\label{exr:viii05-21}
How is \(\mathbb RP^n\) built from \(\mathbb RP^{n-1}\)?
\end{exercise}
\begin{hint}
Use projective filtration.

% VIII-PEDAGOGY-HINT-v1.1 exr:viii05-21 append
The projective chart with last homogeneous coordinate nonzero is an open \(n\)-cell; its complement is the old projective subspace.
\end{hint}
\begin{solution}
Attach one \(n\)-cell.
\end{solution}

\begin{exercise}\label{exr:viii05-22}
What is an \(n\)-skeleton informally?
\end{exercise}
\begin{hint}
Stop at dimension \(n\).

% VIII-PEDAGOGY-HINT-v1.1 exr:viii05-22 append
Check that the boundary of every included cell lands in a lower skeleton, so the stopped construction is closed under attachments.
\end{hint}
\begin{solution}
The union of all cells of dimension at most \(n\).
\end{solution}

\begin{exercise}\label{exr:viii05-23}
Why is attachment a mapping-cone construction?
\end{exercise}
\begin{hint}
A cone on \(S^{n-1}\) is \(D^n\).

% VIII-PEDAGOGY-HINT-v1.1 exr:viii05-23 append
Use a radial parameter to identify the cone on the attaching sphere with a disk, with cone base corresponding to disk boundary.
\end{hint}
\begin{solution}
\(X\cup_\varphi D^n=X\cup_\varphi CS^{n-1}\).
\end{solution}

\begin{exercise}\label{exr:viii05-24}
Why do cell dimensions alone not determine a space?
\end{exercise}
\begin{hint}
Compare degree-\(m\) attachments.

% VIII-PEDAGOGY-HINT-v1.1 exr:viii05-24 append
Compare the constant and identity attachments of a two-cell to a circle: the cell dimensions agree, but one retains a circle homotopy class and the other fills it.
\end{hint}
\begin{solution}
Different attaching maps with identical cell counts can yield different topology.
\end{solution}


\section{Solved problem dossiers}

\begin{problem}[Adjunction universal property]\label{prob:viii05-01}
Prove the gluing criterion for maps out of \(X\cup_\varphi D^n\).
\end{problem}
\begin{solution}
A map out restricts to compatible maps on \(X\) and \(D^n\). Conversely, compatible maps are constant on the quotient equivalence classes and therefore descend uniquely by the quotient property.
\end{solution}

\begin{problem}[One-cell loop]\label{prob:viii05-02}
Attach one interval to one vertex at both endpoints and identify the result.
\end{problem}
\begin{solution}
The quotient \(D^1/\partial D^1\) is a circle.
\end{solution}

\begin{problem}[One-cell edge]\label{prob:viii05-03}
Attach an interval between two distinct vertices.
\end{problem}
\begin{solution}
The result is a graph with two vertices joined by one edge, homeomorphic to a closed interval.
\end{solution}

\begin{problem}[Parallel edges]\label{prob:viii05-04}
Attach two one-cells between the same two vertices. Determine the homotopy type.
\end{problem}
\begin{solution}
The two edges form a circle: go out along one and return along the other.  The space is homeomorphic to \(S^1\).
\end{solution}

\begin{problem}[Bouquet of circles]\label{prob:viii05-05}
Attach \(r\) one-cells with both endpoints at one vertex.
\end{problem}
\begin{solution}
Each interval becomes a circle and all circles meet only at the common vertex, giving \(\bigvee_{j=1}^rS^1\).
\end{solution}

\begin{problem}[Constant attachment]\label{prob:viii05-06}
Prove \(D^n/S^{n-1}\cong S^n\).
\end{problem}
\begin{solution}
The quotient is the one-point compactification of the open disk \(\mathring D^n\cong\mathbb R^n\), hence \(S^n\).
\end{solution}

\begin{problem}[Minimal sphere]\label{prob:viii05-07}
Give the two-cell structure of \(S^n\).
\end{problem}
\begin{solution}
Start from one vertex and attach one \(n\)-disk by the constant boundary map.  The quotient is \(D^n/S^{n-1}\cong S^n\).
\end{solution}

\begin{problem}[Degree-one two-cell]\label{prob:viii05-08}
Identify \(S^1\cup_{\operatorname{id}}D^2\).
\end{problem}
\begin{solution}
The existing circle becomes exactly the boundary of the new disk, so the result is \(D^2\).
\end{solution}

\begin{problem}[Constant two-cell]\label{prob:viii05-09}
Identify \(S^1\cup_cD^2\).
\end{problem}
\begin{solution}
The new disk boundary collapses at the chosen point, turning the disk into \(S^2\) wedged with the old \(S^1\).
\end{solution}

\begin{problem}[Degree-two two-cell]\label{prob:viii05-10}
Explain the \(\mathbb RP^2\) disk model.
\end{problem}
\begin{solution}
Identify antipodal points of \(\partial D^2\).  The quotient boundary is \(\mathbb RP^1\cong S^1\), and its attaching map is the two-to-one map of degree \(2\).
\end{solution}

\begin{problem}[Degree-\(m\) family]\label{prob:viii05-11}
Why can \(X_m=S^1\cup_{\varphi_m}D^2\) vary with \(m\) despite fixed cell counts?
\end{problem}
\begin{solution}
The boundary winds \(m\) times around the one-skeleton.  Cellular homology later records this as multiplication by \(m\), so the attaching data are topologically significant.
\end{solution}

\begin{problem}[Torus polygon]\label{prob:viii05-12}
Derive the torus attaching word.
\end{problem}
\begin{solution}
Label the square boundary in order \(a,b,a^{-1},b^{-1}\); after opposite edge identifications this becomes the loop \(aba^{-1}b^{-1}\).
\end{solution}

\begin{problem}[Genus-two surface]\label{prob:viii05-13}
Give a cell attachment for the closed genus-two surface.
\end{problem}
\begin{solution}
Use one vertex, four one-cells \(a_1,b_1,a_2,b_2\), and attach one two-cell along \([a_1,b_1][a_2,b_2]\).
\end{solution}

\begin{problem}[Projective filtration]\label{prob:viii05-14}
Why does \(\mathbb RP^n\) have one cell in every dimension?
\end{problem}
\begin{solution}
Each difference \(\mathbb RP^k\setminus\mathbb RP^{k-1}\) is \(\mathbb R^k\), so stage \(k\) adds exactly one open \(k\)-cell.
\end{solution}

\begin{problem}[Simultaneous cells]\label{prob:viii05-15}
Describe the quotient for a family of \(n\)-cell attachments.
\end{problem}
\begin{solution}
Use \(X\sqcup\coprod_\alpha D^n_\alpha\) and identify each boundary point \(u\) with \(\varphi_\alpha(u)\); disk interiors remain mutually disjoint.
\end{solution}

\begin{problem}[Characteristic boundary]\label{prob:viii05-16}
Why can a characteristic map fail to be injective on the boundary?
\end{problem}
\begin{solution}
If distinct \(u,v\in S^{n-1}\) satisfy \(\varphi(u)=\varphi(v)\), the quotient identifies them.
\end{solution}

\begin{problem}[Open-cell embedding]\label{prob:viii05-17}
Why is the interior embedded?
\end{problem}
\begin{solution}
The attachment relation only touches boundary points.  No two disk interior points are identified, and none is identified with the old space.
\end{solution}

\begin{problem}[Pushout uniqueness]\label{prob:viii05-18}
State the uniqueness clause for the cell-attachment pushout.
\end{problem}
\begin{solution}
Compatible maps on \(X\) and \(D^n\) determine exactly one map on the quotient because their images cover the adjunction space.
\end{solution}

\begin{problem}[Successive attachments]\label{prob:viii05-19}
What data determine a two-dimensional complex built on a graph?
\end{problem}
\begin{solution}
The one-skeleton graph plus one attaching loop \(S^1\to X^1\) for every two-cell.
\end{solution}

\begin{problem}[Attaching word]\label{prob:viii05-20}
Why can a word in graph edges specify a two-cell attachment?
\end{problem}
\begin{solution}
Reading the oriented edges in sequence defines a loop in the graph, hence a map \(S^1\to X^1\) used as the attaching map.
\end{solution}

\begin{problem}[Cell as cone]\label{prob:viii05-21}
Show an \(n\)-cell attachment is a mapping cone of a sphere map.
\end{problem}
\begin{solution}
\(CS^{n-1}\cong D^n\), with the cone base equal to the disk boundary, so \(X\cup_\varphi CS^{n-1}\) is the same adjunction space.
\end{solution}

\begin{problem}[Homotopic attaching maps boundary]\label{prob:viii05-22}
What conclusion is expected for homotopic attaching maps, and why is it reserved for VIII/08?
\end{problem}
\begin{solution}
Under standard hypotheses their adjunction spaces are homotopy equivalent, but the proof requires homotopy extension/pushout control; VIII/08 supplies that theorem rather than treating the quotients as automatically homeomorphic.
\end{solution}

\begin{problem}[Quotient topology]\label{prob:viii05-23}
Why is the quotient topology essential?
\end{problem}
\begin{solution}
It is what makes the gluing universal property valid and determines continuity of maps into and out of the adjunction space.
\end{solution}

\begin{problem}[Bridge to CW complexes]\label{prob:viii05-24}
How does repeated cell attachment lead to VIII/06?
\end{problem}
\begin{solution}
Build \(X^0,X^1,X^2,\ldots\) dimension by dimension; a CW complex formalizes this filtration together with closure-finiteness and weak-topology conditions.
\end{solution}
